\documentclass[10pt,a4paper]{amsart}
\usepackage[margin=19mm]{geometry}
\usepackage{amsmath,amssymb,mathtools}
\usepackage[scr=rsfs,cal=euler]{mathalfa}
\usepackage{xfrac}
\usepackage[unicode,hidelinks,bookmarks=false]{hyperref}
\newcommand{\ie}{i.e.\ }
\newcommand{\cf}{cf.\ }
\newcommand{\resp}{respectively\ }
\newcommand{\Subsection}{\subsection}
\providecommand{\nobreakdash}{\nobreak\hbox{-}}
\setlength{\emergencystretch}{2em}
\multlinegap0pt
\usepackage{amssymb}
\newtheorem{thm}{Theorem}
\newcommand{\R}{\mathbb{R}}
\title{On the entropy for indeterminate moment problems}
\author{Christian Berg}
\date{Source: arXiv:2511.12684v2 (CC BY 4.0)}
\begin{document}
\maketitle

\noindent\emph{Source and licence.} On the entropy for indeterminate moment problems. Version arXiv:2511.12684v2 (CC BY 4.0), distributed by arXiv under the Creative Commons Attribution 4.0 International licence, http://creativecommons.org/licenses/by/4.0/, as recorded in the arXiv licence metadata for this exact version and shown on the abstract page https://arxiv.org/abs/2511.12684v2. Full licence notice: \texttt{licenses/arxiv-2511.12684-CC-BY-4.0.txt}.\par\smallskip

\noindent\textit{Editorial scope.} Counted unit: the article's thm thm:bd (source0.tex lines 316--317). The article's complete original proof of it is delivered with it (source0.tex lines 319--375, one \texttt{proof} environment of 57 lines). No mathematics is written, repaired or re-derived: the statement and the proof are the article's own lines, in the article's own notation, and no retained line is substituted or re-flowed.  No further source-local context is needed: every cross-reference of every retained line resolves inside the two delivered spans.  Nothing was omitted from any retained span, so the package carries no omission record.  No retained line contains a citation, so the wrapper carries no bibliography and no bibliography entry.  No retained line contains a figure environment, an image-inclusion command or drawing code, so no image is delivered and the package declares no asset dependency.  Editorial formatting only, in addition: an AMS \texttt{amsart} preamble that restates the article's own declarations and macros used by the retained lines, namely the environments \texttt{thm} on separate counters, together with the standard packages those lines need.  The retained environments print in this document as Theorem 1 in order of appearance, whereas the article prints the same items with the numbering of its own full document.  The complete native source of the article as distributed in the arXiv e-print for this exact version is bundled unchanged as \texttt{sources/189220/source0.tex}.
\begin{thm}\label{thm:bd} Assume $1<a<1/q$. The function $\varphi(x):=[x]_\infty^2+[x/a]_\infty^2$ is bounded below on $\R$ by a positive constant $LB(a,q)$  given in \eqref{eq:LB}.
\end{thm}

\begin{proof} First notice that for $x\le \xi<1$ we have $[x]_\infty \ge [\xi]_\infty>0$, and hence
\begin{equation}\label{eq:1}
\varphi(x)\ge [\xi]_\infty^2,\quad x\le \xi<1.
\end{equation}
 We choose constants $\alpha,\beta$ such that
\begin{equation}
1<\alpha<a<\beta<1/q.
\end{equation} 
We first estimate $[x]_\infty^2$ from below when $x$ belongs to the interval $[\alpha q^{-n},\beta q^{-n}]$ for fixed $n=0,1,\ldots.$

To obtain the estimates we use that the parabola $(1-cx)^2$ has minimum at $x=1/c$, so the minimum of the parabola over an interval $[l,r]$ is achieved at the right endpoint $r$ of the interval when $r<1/c$ and at the left endpoint $l$ when $1/c<l$.
 
For $x\in[\alpha q^{-n}, \beta q^{-n}]$ we find
\begin{eqnarray*}
[x]_\infty^2&=&\left(\prod_{k=0}^{n-1}(1-xq^k)^2\right) (1-xq^n)^2\prod_{k=n+1}^\infty (1-xq^k)
^2\\
&\ge& \left(\prod_{k=0}^{n-1}(1-\alpha q^{k-n})^2\right)(1-\alpha)^2\prod_{k=n+1}^\infty(1-\beta q^{k-n})^2\\
&=& (\alpha-1)^2\left(\prod_{j=1}^n\left(\frac{\alpha}{q^j}\right)^2(1-q^j/\alpha)^2\right)\prod_{j=1}^\infty(1-\beta q^j)^2\\
&=& (\alpha-1)^2\alpha^{2n}q^{-n(n+1)}[q/\alpha]_n^2\,[\beta q]_\infty^2\\
&>&(\alpha-1)^2\alpha^{2n}q^{-n(n+1)}[q/\alpha]_\infty^2\,[\beta q]_\infty^2=:K_n.
\end{eqnarray*}
Note that $K_n$ is increasing in $n$ so $K_n\ge K_0$ and hence
\begin{equation}\label{eq:2}
\varphi(x)\ge K_0,\quad x\in\bigcup_{n=0}^\infty [\alpha q^{-n},\beta q^{-n}].
\end{equation} 

We next estimate $[x/a]_\infty^2$ from below for $x\in [\beta q^{-n+1},\alpha q^{-n}]$
for fixed $n=0,1,\ldots.$

We find
\begin{eqnarray*}
[x/a]_\infty^2&=&\left(\prod_{k=0}^{n-1}(1-(x/a)q^k)^2\right)(1-(x/a)q^n)^2\prod_{k=n+1}^\infty (1-(x/a)q^k)^2\\
&\ge&\left(\prod_{k=0}^{n-1}(1-(q\beta/a)q^{k-n})^2\right)(1-\alpha/a)^2\prod_{k=n+1}^\infty (1-(\alpha/a)q^{k-n})^2\\
&=& (1-\alpha/a)^2\left(\prod_{j=0}^{n-1} \left(\frac{\beta}{aq^j}\right)^{2}(1-(a/\beta)q^j)^2\right)\prod_{j=1}^\infty(1-(\alpha/a)q^j)^2\\
&=&(1-\alpha/a)^2(\beta/a)^{2n}q^{-n(n-1)}[a/\beta]_n^2\,[q\alpha/a]_\infty^2\\
&>&(1-\alpha/a)^2(\beta/a)^{2n}q^{-n(n-1)}[a/\beta]_\infty^2\,[q\alpha/a]_\infty^2=:L_n.
\end{eqnarray*}
Note that $L_n$ is increasing in $n$ so $L_n\ge L_0$ and hence
\begin{equation}\label{eq:3}
\varphi(x)\ge L_0,\quad x\in\bigcup_{n=0}^\infty [\beta q^{1-n},\alpha q^{-n}].
\end{equation}
From \eqref{eq:1} we have $\varphi(x)\ge [q\beta]_\infty^2$ for $x\le q\beta$.
Combining this with \eqref{eq:2} and \eqref{eq:3}, we see that $\varphi(x)$ is bounded below on $\R$ by the constant
$$
\varphi(x)\ge \min\{K_0,L_0,[q\beta]_\infty^2\}.
$$

To get  a constant which does not depend on the chosen values $\alpha,\beta$ we choose
$\alpha:=\sqrt{a}, \beta:=\sqrt{a/q}$ and get the following lower bound for $\varphi(x)$:
$$
\min\{(\sqrt{a}-1)^2[q/\sqrt{a}]_\infty^2[\sqrt{aq}]_\infty^2, (1/\sqrt{a}-1)^2 [q/\sqrt{a}]_\infty^2[\sqrt{aq}]_\infty^2, [\sqrt{aq}]_\infty^2\},
$$
which can be simplified to
\begin{equation}\label{eq:LB}
LB(a,q):=[\sqrt{aq}]_\infty^2[q/\sqrt{a}]_\infty^2\min\{(1/\sqrt{a}-1)^2, 1/[q/\sqrt{a}]_\infty^2\}.
\end{equation}
\end{proof}

\end{document}
